\lesson{From typing derivations to target terms}\label{sec:evidence}
\focus{The translation represents subtyping derivations as evidence and subsumption as a cast.}

The source subsumption rule combines a typing derivation $h$ with a
subtyping derivation $d$:
\[
\src{\frac{h:\Gamma\vdash u:S\qquad d:\Gamma\vdash S<:T}
           {\textsc{Sub}(h,d):\Gamma\vdash u:T}}.
\]
Subsumption changes the assigned type while preserving the source
expression $u$. The chosen conversion is therefore recorded in the
derivation rather than the term.

Write $\tr{h}$ for the target term translated from a source typing
derivation $h$, and $\tr{d}$ for the evidence translated from a source
subtyping derivation $d$. The subsumption clause is
\[
\tgt{\tr{\textsc{Sub}(h,d)}=\cast{\tr{h}}{\tr{d}}.}
\]
This translation clause recursively translates the two premises and
combines their results with the target $\ev{cast}$ constructor.

The target is called \FC{}. An expression $e:S\le T$ is called
\emph{inclusion evidence}, or an \emph{inclusion coercion}.
The target has no implicit subsumption: a cast records a type change
justified by evidence, with no potentially failing runtime type test.
Its typing rule checks the two translated premises:
\[
\tgt{\frac{\Gamma'\vdash \tr{h}:S'\qquad
           \Gamma'\vdash\tr{d}:S'\le T'}
          {\Gamma'\vdash\cast{\tr{h}}{\tr{d}}:T'}}.
\]
Primes denote the translated context and types. We use the same notation
for $W$ and $x.A$ on both sides of the translation.

A target \emph{atom} is a variable, possibly surrounded by typing
wrappers such as casts. Its \emph{root} is the underlying variable:
both $f$ and $\cast{f}{e}$ have root $f$. Atoms serve as the receivers
and arguments of operations.

The \emph{erasure} $\erase{t}$ removes typing annotations and evidence
wrappers, yielding the underlying program. In particular,
$\erase{\cast{u}{e}}=\erase{u}$.
Erasure relates program syntax; the typed target's operational semantics
may still consult evidence.

\subsection*{Translating the two bound selections}
Consider an occurrence of $f$ under the parameter context
$\src{\Gamma=x:\iface{W}{W},\,f:W}$.
This is the client part of the example: for now, assume that its object
parameter already exposes the exact bounds. We will translate this body
first. Section~\ref{sec:closed-program} then supplies the literal's
introduction and recursive opening, displays the complete target program,
and substitutes the allocated atoms into this same body. No assumption
about a stored definition is needed for the present calculation.
Reading the assumed interface gives two source derivations:
\[
\src{\begin{aligned}
d_\ell&:\Gamma\vdash W<:x.A,\\
d_h&:\Gamma\vdash x.A<:W.
\end{aligned}}
\]
Write the translated evidence using the abbreviations
\[
 \ell_x=\ev{lower}_A(x):W\le x.A,\qquad
 h_x=\ev{upper}_A(x):x.A\le W.
\]
Section~\ref{sec:telescopes} expands each abbreviation into an indexed
member read. Figure~\ref{fig:derivation} translates the two subsumption
steps from their premises: variable typing yields the atom $f$ and the
lower-bound derivation yields $\ell_x$, which the inner \textsc{Sub}
combines into $\cast{f}{\ell_x}$. The upper-bound derivation then yields
$h_x$, and the outer \textsc{Sub} adds the second cast.


\begin{figure}[!htbp]
\begin{examplepanel}{srccol}{Source derivation --- two subsumptions on the same occurrence of \(f\)}
\small
\src{\(\Gamma=x:\iface{W}{W},\,f:W\).
Here \(h_f\) derives variable typing for \(f\), and \(d_\ell,d_h\)
derive the lower and upper bounds of \(x.A\).}
\[
\src{
\inferrule*[right=\textsc{Sub}]{
 \inferrule*[right=\textsc{Sub}]{
 h_f:\Gamma\vdash f:W\\
 d_\ell:\Gamma\vdash W<:x.A
 }{\Gamma\vdash f:x.A}\\
 d_h:\Gamma\vdash x.A<:W
}{\Gamma\vdash f:W}}
\]
\end{examplepanel}
\begin{examplepanel}{tgtcol}{Translation of the derivation}
\[
\tgt{
\begin{aligned}
 \tr{h_f}&=f,\\
 \tr{d_\ell}&=\ell_x=\ev{lower}_A(x):W\le x.A,\\
 \tr{d_h}&=h_x=\ev{upper}_A(x):x.A\le W,\\
 \tr{\textsc{Sub}(h_f,d_\ell)}&=\cast{f}{\ell_x},\\
 \tr{\textsc{Sub}(\textsc{Sub}(h_f,d_\ell),d_h)}
   &=\cast{\cast{f}{\ell_x}}{h_x}.
\end{aligned}}
\]
\end{examplepanel}
\begin{examplepanel}{runcol}{Erasure of the target term}
\[
\run{\erase{\cast{\cast{f}{\ell_x}}{h_x}}=f.}
\]
\end{examplepanel}
\caption{Each \textsc{Sub} node contributes a target \(\ev{cast}\),
while the source expression remains \(f\) throughout the derivation.
The calculation assumes parameter binders for \(x\) and \(f\).
Figure~\ref{fig:let-translation} extends it to the program's let bindings.}
\label{fig:derivation}
\end{figure}



\subsection*{Dependence on the source derivation}
The judgment $\src{\Gamma\vdash f:W}$ admits several derivations.
Variable typing translates to $f$, while the two subsumptions translate
to $\cast{\cast{f}{\ell_x}}{h_x}$. A third derivation first composes
the bounds:
\[
\src{\frac{d_\ell:\Gamma\vdash W<:x.A\qquad
           d_h:\Gamma\vdash x.A<:W}
          {\textsc{Trans}(d_\ell,d_h):\Gamma\vdash W<:W}}.
\]
Source transitivity becomes an evidence constructor:
\[
\tgt{\tr{\textsc{Trans}(d_\ell,d_h)}
       =\ev{trans}(\ell_x,h_x).}
\]
One subsequent subsumption therefore gives
$\cast{f}{\ev{trans}(\ell_x,h_x)}$.
The three target expressions are syntactically distinct, although all
erase to $f$. In particular, nested casts and a cast by composed evidence
are different constructor expressions. The translation is thus a
function of the source derivation, not only of its concluding term.

We abbreviate $\ev{trans}(p,q)$ by $p;q$, in derivation order:
\[
 \frac{p:S\le M\qquad q:M\le T}{p;q:S\le T}.
\]
Composition requires matching intermediate endpoints. Here the chain is
$W\xrightarrow{\ell_x}x.A\xrightarrow{h_x}W$.

\subsection*{Let translation and evaluation}
The translation preserves the running program's let bindings:
\[
\tgt{
\tr{\textsc{Let}(h_1,h_2)}
=\ev{let}\ y=\tr{h_1}\ \ev{in}\ \tr{h_2}.}
\]
Figure~\ref{fig:let-translation} applies this clause to the body of the
running example. The first right-hand side uses the lower-bound cast.
The second uses the upper-bound cast on the newly bound name
$\ty{hidden}$. The application also needs the argument cast
$\ev{top}(W):W\le\top$.

Application translates the variable typing derivations of its operands
to atoms. Subsumption in such a derivation contributes a cast wrapper,
so the source argument $f$ becomes $\cast{f}{\ev{top}(W)}$.
The full language distinguishes casts of terms, values, and atoms;
we use the same notation when the syntactic position determines the sort.


\begin{figure}[!htbp]
\begin{examplepanel}{srccol}{Source body --- after binding \(x\) and \(f\)}
\begin{lstlisting}[style=source]
let hidden    : x.A = f in
let recovered : W   = hidden in
recovered f
\end{lstlisting}
\end{examplepanel}
\begin{examplepanel}{tgtcol}{Target translation with explicit casts}
\tgt{\(\ell_x=\ev{lower}_A(x)\), \(h_x=\ev{upper}_A(x)\),
\(t_W=\ev{top}(W):W\le\top\).}
\begin{lstlisting}[style=target,mathescape=true]
let hidden    = cast(f, $\ell_x$) in
let recovered = cast(hidden, $h_x$) in
recovered (cast(f, $t_W$))
\end{lstlisting}
\end{examplepanel}
\begin{examplepanel}{tgtcol}{Evaluation by substitution of the aliases}
\small
\[
\tgt{\begin{aligned}
 &\ev{let}\ \ty{recovered}=\cast{\cast{f}{\ell_x}}{h_x}\ \ev{in}\
       \ty{recovered}\ \cast{f}{t_W}\\
 &\hspace{5em}\longrightarrow^{*}
       \cast{\cast{f}{\ell_x}}{h_x}\ \cast{f}{t_W}.
\end{aligned}}
\]
The first binding substitutes \(\cast{f}{\ell_x}\) for \(\ty{hidden}\);
the second substitutes the nested cast for \(\ty{recovered}\).
The relation \(\longrightarrow^{*}\) includes the administrative steps
that rearrange pending casts.
\end{examplepanel}
\caption{Translation preserves the source let structure; evaluation
produces nested casts by substituting the let-bound atoms. The resulting
application erases to \(f\,f\). As in Figure~\ref{fig:derivation}, the
exact interface is assumed directly.
Section~\ref{sec:closed-program} adds the literal's introduction and
recursive opening, then instantiates its binders with atoms over the store.}
\label{fig:let-translation}
\end{figure}



During evaluation, substitution of the let-bound atoms produces the
nested wrappers. This display still has the parameters $x$ and $f$ free.
Closing them is a further calculation, given in
Section~\ref{sec:closed-program}; it must justify both the object
interface used by the evidence and the function supplied to the call.

\subsection*{Checking evidence endpoints}
The expression $\ell_x;\ev{top}(W)$ fails the composition rule:
its first component ends at $x.A$ and its second starts at $W$.
Inserting the upper-bound evidence gives the well-typed composition
$\ell_x;h_x;\ev{top}(W)$. Alternatively,
$\ell_x;\ev{top}(x.A)$ reaches $\top$ directly from $x.A$.

Even if a later store resolves $x.A$ to $W$, checking still compares
syntactic endpoints. Store-based normalization establishes what the
evidence means at runtime.

\subsection*{Erasure and operational correspondence}
The erasure equation $\erase{\cast{u}{e}}=\erase{u}$ helps show that
translation preserves the underlying erased program.
The typed target machine still consults a function atom's
evidence to perform a call. A simulation theorem relates those typed
steps to the erased program's execution.

\remember{Translation constructs evidence from source derivations;
evaluation substitutes atoms and consults their evidence; erasure
removes the evidence to relate the underlying programs.}

\exercise{order}{What does $h_x;\ell_x$ prove? Is it the same expression
as $\ell_x;h_x$? Also give explicit evidence for $W\le\top$ that goes
through $x.A$.}

\exercise{translate}{The source occurrence $f:W$ is first used at
$x.A$, then at $\top$ by the source \textsc{Top} rule $x.A<:\top$.
Give its target term when there are two \textsc{Sub} nodes, and when
there is one \textsc{Sub} over a \textsc{Trans} node. What does each erase to?}
